% P05.2 one-dimensional graph classification proof, 25 September 2026.

\begin{proof}[\Cref{lem:frontier-polyhedral-ambient-vertex-accessibility}]
Choose a rational polyhedral norm $\lVert\cdot\rVert$ on $N_{\mathbf R}$.
On every unbounded edge $e$, let $r_e>0$ be the eventual slope of $\gamma$
in the primitive outgoing direction $u_e$.  Since there are only finitely
many such edges, choose $\varepsilon>0$ such that
\[
                         \varepsilon\lVert u_e\rVert<r_e
                         \qquad\text{for every unbounded }e.
\]
If $\mathbf X$ has no unbounded edge, any positive $\varepsilon$ will do.

Fix a vertex $v$ and choose $C>\gamma(v)$.  The restriction to $\mathbf X$
of the ambient convex function
\[
                         h(x):=C+\varepsilon\lVert x-v\rVert
\]
is polyhedrally convex and is affine on a finite rational polyhedral
refinement.  A finite maximum preserves every polyhedral Jensen inequality,
so the continuous function
\[
                         F:=\max\{\gamma,h\}
\]
is polyhedrally convex and has only finitely many affine pieces on every
edge.  Its breakpoints need not be rational, so one further rational
interpolation is required.

Along every unbounded edge the eventual slope of $\gamma-h$ is positive.
Hence $F=\gamma$ outside a compact subset of $\mathbf X$.  On each edge,
choose finitely many rational points which include its endpoints, meet the
affine germ of $F$ at every incident vertex, and, on an unbounded edge,
include a point beyond which $F=\gamma$.  Let $\varphi$ be the affine
interpolation of the values of $F$ at these points, extended by $\gamma$ on
the remaining unbounded tails.  Secant slopes of a convex function on an
interval are nondecreasing, so $\varphi$ is convex on each edge.  Its
outgoing slopes at every original vertex equal those of $F$; hence all local
nonnegative-cycle inequalities are preserved.  The one-dimensional local
slope criterion gives
\[
 \varphi\in\mathrm{PA}_{\mathbf R}(\mathbf X)\cap
 \mathrm{PC}(\mathbf X),
 \qquad
 \varphi-\gamma\in\mathrm{PA}_{\mathbf R,b}(\mathbf X).
\]

The strict inequality $C>\gamma(v)$ persists on a sufficiently small closed
star about $v$.  Choose the first rational interpolation points inside this
star.  Then $\varphi=h$ on their incident segments, and its outgoing slope
along every incident edge is
$\varepsilon\lVert u_{e/v}\rVert>0$.  Therefore
\[
 \MA_{\mathrm{poly},\mathbf R}(\varphi)\{v\}
 =\varepsilon\sum_{e\in E(v)}w_e\lVert u_{e/v}\rVert>0.
\]
Thus $v$ is $\gamma$-mass-accessible.
\end{proof}

\begin{proof}[\Cref{lem:frontier-polyhedral-accessible-mass-obstruction}]
Choose $\varphi$ as in
\Cref{def:frontier-polyhedral-mass-accessible}.  Let
$s=(s_e)_{e\in E(v)}$ be its outgoing slope vector at $v$, let
$w=(w_e)_{e\in E(v)}$, and put
\[
                         m:=w\mathbin{\cdot}s>0.
\]
The local convexity criterion gives
$\lambda\mathbin{\cdot}s\geq0$ for every $\lambda\in K_v$.

Let $L_v$ be the linear span of $K_v$.  Since $\dim K_v>1$ and
$w\mathbin{\cdot}s>0$, choose a positive rational vector
$c=(c_e)_{e\in E(v)}$ such that the restrictions of $c$ and $s$ to $L_v$
are not proportional.  Multiplying $c$ by a sufficiently large positive
rational number, choose points $x_e$ at integral-affine distance
$\rho_e=c_e^{-1}$ from $v$ such that every segment $[v,x_e]$ contains no
other vertex and $\varphi$ is affine on it.  The slice
\[
                       B_c:=\{\lambda\in K_v:\lambda\mathbin{\cdot}c=1\}
\]
is nonempty and compact.  Set
\[
                       \tau:=\min_{\lambda\in B_c}
                       \lambda\mathbin{\cdot}s\geq0.
\]

Raise $\varphi(v)$ by $\tau$, interpolate affinely to the unchanged values
$\varphi(x_e)$, and leave $\varphi$ unchanged off the closed star.  Denote
the resulting function by $\widetilde\varphi$.  Its outgoing slope vector
at $v$ is
\[
                         \widetilde s=s-\tau c.
\]
The definition of $\tau$ gives
$\lambda\mathbin{\cdot}\widetilde s\geq0$ for every $\lambda\in K_v$,
and a minimizer $\lambda^0\in B_c$ satisfies
\begin{equation}
                  \lambda^0\mathbin{\cdot}\widetilde s=0.
                  \label{eq:frontier-polyhedral-supporting-relation}
\end{equation}
At each $x_e$ the modified slope increases back to the old slope, so
one-variable convexity is preserved.  All other local inequalities are
unchanged.  Thus $\widetilde\varphi$ is polyhedrally convex and remains
$\gamma$-bounded.

Its mass at $v$ is still positive.  This is immediate when $\tau=0$.
If $\tau>0$, the linear functional
$\lambda\mapsto\lambda\mathbin{\cdot}s$ is nonconstant on $B_c$ by the
choice of $c$.  The point
$w/(w\mathbin{\cdot}c)$ lies in the relative interior of $B_c$, so its
value is strictly larger than the minimum.  Hence
\begin{equation}
 \frac{w\mathbin{\cdot}s}{w\mathbin{\cdot}c}>\tau,
 \qquad
 w\mathbin{\cdot}\widetilde s>0.
 \label{eq:frontier-polyhedral-positive-projected-mass}
\end{equation}

Let $b$ be the real-piecewise-affine function which is one at $v$, zero at
every $x_e$, affine on $[v,x_e]$, and zero off the closed star.  Put
\[
                         a:=\widetilde\varphi+b,
                         \qquad
                         u:=a-\gamma.
\]
Then $u$ is bounded and real-piecewise affine, and
$\widetilde\varphi$ is an admissible competitor for $P_\gamma(u)$.

Let $q\leq a$ be any polyhedrally convex envelope competitor and put
\[
 C_0:=\sum_e\frac{\lambda^0_e}{\rho_e},
 \qquad
 \alpha_e:=\frac{\lambda^0_e/\rho_e}{C_0}.
\]
Because $x_e=v+\rho_eu_{e/v}$ and $\lambda^0\in K_v$, the points $x_e$
with coefficients $\alpha_e$ form a polyhedral convex combination with
barycentre $v$.  Since $b(x_e)=0$, polyhedral convexity and
\eqref{eq:frontier-polyhedral-supporting-relation} give
\[
 \begin{aligned}
 q(v)
 &\leq\sum_e\alpha_eq(x_e)
 \leq\sum_e\alpha_e\widetilde\varphi(x_e)\\
 &=\widetilde\varphi(v)
   +\frac{1}{C_0}\sum_e\lambda^0_e\widetilde s_e
 =\widetilde\varphi(v).
 \end{aligned}
\]
Taking the supremum over competitors and using the competitor
$\widetilde\varphi$ shows that, for $P:=P_\gamma(u)$,
\[
                         P(v)=\widetilde\varphi(v),
                         \qquad
                         P\geq\widetilde\varphi.
\]

The one-dimensional real-affine envelope theorem makes $P$
real-piecewise affine.  Since $P-\widetilde\varphi$ is nonnegative and
vanishes at $v$, every outgoing slope of $P$ at $v$ is at least the
corresponding slope of $\widetilde\varphi$.  By
\eqref{eq:frontier-polyhedral-positive-projected-mass},
\[
              \MA_{\mathrm{poly},\mathbf R}(P)\{v\}
              \geq w\mathbin{\cdot}\widetilde s>0.
\]
At $v$ one has $P-a=-1$, while $P-a\leq0$ everywhere and
$\MA_{\mathrm{poly},\mathbf R}(P)$ is nonnegative.  Therefore
\[
 \int_{\mathbf X}(P-a)\MA_{\mathrm{poly},\mathbf R}(P)
 \leq-\MA_{\mathrm{poly},\mathbf R}(P)\{v\}<0.
\]
This is the required failure of orthogonality.
\end{proof}

\begin{proof}[\Cref{thm:frontier-polyhedral-graph-orthogonality-classification}]
If $\dim K_v=1$ at every vertex, orthogonality is
\Cref{thm:frontier-polyhedral-local-rank-one-graph}.

Conversely, suppose that $\dim K_v>1$ at some vertex.  By
\Cref{lem:frontier-polyhedral-ambient-vertex-accessibility}, this vertex is
$\gamma$-mass-accessible.  Then
\Cref{lem:frontier-polyhedral-accessible-mass-obstruction} gives a bounded
real-piecewise-affine obstacle with a strictly negative orthogonality
integral.  Thus orthogonality fails.
\end{proof}
